A diffusion model is a denoiser trained at every noise level. When the data are first-order Markov sequences, such as
time series or one-dimensional physical fields, we ask what a denoiser gains from this structure and how it should use
it. The structure is not inherited by the score: for a Gaussian chain the precision matrix of the noisy data, tridiagonal
for clean data, fills in at intermediate noise. It survives in the posterior. A Markov chain observed through independent
noise is a tree, so belief propagation (BP) gives the exact posterior mean, and hence the exact score, at every noise
level once the transition kernel is known. Learning the denoiser thus reduces to learning a function of two variables,
which we do by expectation--maximisation with BP as the E-step, from noisy sequences alone and without differentiating
through BP. On a chain with Laplace innovations, kernels learned in this way give scores 2.1 to 2.8 times more accurate
than the best of six neural denoisers trained on the same data. A neural energy kernel, which assumes nothing beyond the
Markov property, matches the best hand-designed family and beats every network in at least 15 of 16 seeds at every
training-set size. On nine other chains, kernel families with a fixed innovation law fall to the error of a linear
denoiser once the chain leaves their class, while the neural kernel does not. The advantage shrinks as the data depart
from the Markov property, at a rate set by their distance to the closest first-order chain, and BP on pairs of sites
restores it for a second-order chain. On the stochastic Allen--Cahn field, the neural kernel, with at most 1{,}185
parameters, beats a diffusion U-Net with $1.2\times10^{7}$ parameters in 30 of 32 paired seeds at 2{,}048 training
sequences, and a five-parameter kernel with the form of the lattice action recovers the temperature to within one
percent.
% src: 2.1-2.8: reports/paper-figures/README.md (E02a em_bp vs E03 networks, paired, 16 seeds, N=32..2048)
% src: neural energy kernel: tools/core/read_e56.py (E56 frozen 1 Nov 2026): energy 0.0235 vs em_bp 0.0231 at N=2048;
%      paired wins against each of E03's six networks 16/16 at every N except BiGRU at N=512 (15/16)
% src: nine chains: figures/tab_chains.tex (tools/paper_figs/make_paper_figures.py::tab_chains on frozen E57 and E61)
% src: second-order chain: tools/core/read_e62.py (E62 frozen 6 Nov 2026)
% src: U-Net: E55 finished cells (frozen 7 Nov 2026, not yet read against the predictions), seeds 0-7, N=2048: energy below U-Net 8+8+8+6 of 32;
%      n_params em_energy 337 (1,185 at Theta=0.5), unet 11,774,977
% src: temperature: E33 (Theta-hat/Theta 0.9995-1.0068, answers note Q8)
